

\documentclass[12pt,a4paper]{report}

% Required Packages
\usepackage[utf8]{inputenc}
\usepackage[T1]{fontenc}
\usepackage{geometry}
\usepackage{setspace}
\usepackage{graphicx}
\usepackage{float}
\usepackage{amsmath}
\usepackage{booktabs}
\usepackage[hidelinks]{hyperref}
\usepackage{ragged2e}
% Page Formatting
\geometry{
    left=1.25in,
    right=1in,
    top=1in,
    bottom=1in
}
\onehalfspacing

\begin{document}

% =====================================
% TITLE PAGE
% =====================================


% =====================================
% TITLE PAGE
% =====================================

\begin{titlepage}
    \centering

    \vspace*{0.5cm}

    % College Logo
    % Uncomment after uploading college_logo.png
    % \includegraphics[width=3.5cm]{college_logo.png}

    \vspace{0.8cm}

    {\fontsize{20}{24}\selectfont
    \textbf{FUNDAMENTALS OF DATA SCIENCE}\par}

    \vspace{1.2cm}

    {\fontsize{34}{40}\selectfont
    \textbf{SALES INSIGHTS}\par}

    \vspace{0.5cm}

    {\fontsize{18}{23}\selectfont
    \textit{Data Analysis and Machine Learning}\par}

    \vspace{1.3cm}

    {\fontsize{24}{29}\selectfont
    \textbf{PROJECT REPORT}\par}

    \vspace{0.7cm}

    {\fontsize{18}{23}\selectfont
    \textbf{Fundamentals of Data Science}\par}

    \vspace{1.3cm}

    {\fontsize{20}{25}\selectfont
    \textbf{Submitted By}\par}

    \vspace{0.6cm}

    {\fontsize{16}{22}\selectfont
    \begin{tabular}{ll}
        \textbf{Student Name} &
        \textbf{SAP ID} \\[0.35cm]

        Arnav Rai & 590027562 \\[0.25cm]
        Ishi Agarwal & 590025685 \\[0.25cm]
        Vasvi Sood & 590027632
    \end{tabular}
    }

    \vfill

    {\fontsize{18}{23}\selectfont
    \textbf{B.Tech -- Semester III}\par}

    \vspace{0.5cm}

    {\fontsize{17}{22}\selectfont
    Academic Year: 2026--2027\par}

\end{titlepage}


\pagenumbering{roman}

\tableofcontents
\newpage

\pagenumbering{arabic}

% =====================================
% 1. INTRODUCTION
% =====================================

\chapter{Introduction}

\begingroup
\fontsize{14}{21}\selectfont

The primary purpose of the \textbf{Sales Insights}
project is to apply the concepts learned in our
Fundamentals of Data Science course to a sales
transaction dataset. The dataset consists of
100 records and 58 attributes, including
customer information, product details, sales
amounts, profits, and preprocessed features.

Through this project, we have worked with
data preprocessing, statistical hypothesis
testing using the \textbf{Z-test}, and
dimensionality reduction using
\textbf{Principal Component Analysis (PCA)}.

Our main objective is to understand the
practical implementation of these techniques
and connect the theoretical concepts studied
in class with their application to a dataset.

\endgroup

\chapter{Problem Statement}

\begingroup
\fontsize{14}{21}\selectfont
\justifying

In the modern business environment, organizations generate
large volumes of sales transaction data through their
daily operations. This data contains valuable information
regarding customer purchasing behaviour, product sales,
regional demand, pricing strategies, discounts, and
overall business performance. However, raw sales data
often contains missing values, inconsistencies, duplicate
records, and variations in data formats, making it
difficult to extract meaningful and reliable insights
directly.

The primary problem addressed in this project,
\textbf{Sales Insights}, is the difficulty of transforming
raw sales transaction data into structured, meaningful,
and actionable information that can support effective
business decision-making. Without systematic data
analysis, businesses may struggle to identify important
sales trends, understand customer preferences, evaluate
the impact of discounts, and determine the factors
influencing sales performance.

Furthermore, traditional analysis methods may not
adequately capture complex relationships between
multiple variables present in sales datasets.
Factors such as customer age, marketing expenditure,
discount percentage, customer ratings, and delivery
duration may influence business outcomes in different
ways. Therefore, a systematic data science approach
is required to examine these relationships and
identify significant patterns.

To address these challenges, this project applies
\textbf{Data Preprocessing}, \textbf{Exploratory Data
Analysis (EDA)}, \textbf{Feature Engineering},
\textbf{Principal Component Analysis (PCA)},
\textbf{Linear Regression}, and
\textbf{Statistical Hypothesis Testing}.

Data preprocessing is used to improve data quality
by handling missing values, duplicate records,
and other inconsistencies. Exploratory Data Analysis
helps identify trends, distributions, and relationships
between different attributes. Feature engineering
prepares relevant variables for further analysis,
while Principal Component Analysis assists in
understanding the underlying structure of
multidimensional data.

Linear Regression is employed to investigate
relationships between selected independent
variables and a continuous target variable.
Statistical hypothesis testing provides a
systematic framework for evaluating assumptions
and determining whether observed differences
are statistically significant.

The overall aim of the project is to develop
a structured analytical approach that converts
sales transaction data into useful business
insights, supports evidence-based decision-making,
and demonstrates the practical application of
fundamental data science techniques.

\endgroup


\chapter{Project Objectives}

\begingroup
\fontsize{14}{21}\selectfont

The primary objective of the \textbf{Sales Insights}
project is to analyse sales transaction data using
statistical methods, data visualization, and machine
learning techniques to identify meaningful patterns,
understand customer purchasing behaviour, and
support data-driven business decision-making.

The specific objectives of this project are as follows:

\begin{enumerate}

    \item \textbf{Product Category Analysis:}
    To identify which product categories generate
    the highest revenue and evaluate their
    contribution to overall sales performance.

    \item \textbf{Regional Sales Analysis:}
    To determine which geographical regions
    contribute the most to total sales and
    compare sales performance across different
    locations.

    \item \textbf{Customer Behaviour Analysis:}
    To examine how customer characteristics,
    including age, customer type, and purchasing
    preferences, influence buying behaviour
    and sales patterns.

    \item \textbf{Discount and Profitability Analysis:}
    To investigate the relationship between
    discount percentages, sales revenue, and
    business profitability.

    \item \textbf{Identification of Sales Factors:}
    To identify the important factors associated
    with higher sales, including marketing
    expenditure, discounts, and customer-related
    attributes.

    \item \textbf{Sales Prediction Using Linear Regression:}
    To develop a Linear Regression model for
    predicting sales using relevant independent
    variables and evaluate its predictive
    performance.

    \item \textbf{Dimensionality Reduction Using PCA:}
    To apply Principal Component Analysis (PCA)
    to reduce dataset dimensionality while
    retaining a substantial proportion of
    the original variance.

    \item \textbf{Statistical Hypothesis Testing:}
    To determine whether statistically significant
    differences in sales exist between selected
    customer segments using an appropriate
    statistical hypothesis test.

\end{enumerate}

Through these objectives, the project aims to
demonstrate the practical application of
Fundamentals of Data Science concepts in
transforming raw sales data into meaningful
business insights.

\endgroup


% =====================================
% CHAPTER 4: DATASET DESCRIPTION
% =====================================

\chapter{Dataset Description}

\begingroup
\fontsize{14}{21}\selectfont

The \textbf{Sales Insights Dataset} is a structured
dataset containing synthetically generated sales
transaction records. It is used for practising
data preprocessing, statistical analysis, and
dimensionality reduction techniques.

The original dataset contains 100 transactions
and 21 attributes. After preprocessing and
feature engineering, the dataset contains
100 records and 58 columns.

% =====================================
% 4.1 DATASET SOURCE
% =====================================

\section{Dataset Source}

The dataset consists of synthetic sales
transaction data prepared for educational
and analytical purposes. It does not contain
actual business transactions.

The transactions cover the period from
\textbf{16 January 2025 to 23 December 2025}
and represent sales across 12 cities
belonging to four geographical regions
of India: North, South, East, and West.

The dataset contains information related
to customers, products, transaction values,
discounts, marketing expenditure, and profit.
Monetary values are represented in
Indian Rupees (INR).

Since the data is artificially generated,
the findings obtained from the analysis
are applicable to the given dataset and
should not be generalized to real businesses.

% =====================================
% 4.2 DATASET ATTRIBUTES
% =====================================

\section{Dataset Attributes}

The original dataset consists of
\textbf{100 rows and 21 columns},
representing 2,100 data cells.

Each row represents a sales transaction
identified by a unique
\texttt{Transaction\_ID}.
The dataset contains 50 distinct
customer identifiers, indicating that
some customers appear in multiple transactions.

The attributes provide information
about customers, products, sales,
and transaction-related activities.

The important attributes include:

\begin{itemize}

    \item \textbf{Transaction\_ID:}
    Unique identifier assigned to each
    sales transaction.

    \item \textbf{Date:}
    Date on which the transaction occurred.

    \item \textbf{Customer\_ID and Customer\_Name:}
    Attributes used to identify customers.

    \item \textbf{Customer\_Age and Customer\_Type:}
    Customer-related information, including
    age and customer classification.

    \item \textbf{City and Region:}
    Geographical information associated
    with each transaction.

    \item \textbf{Product\_ID and Product\_Category:}
    Information identifying the product
    and its category.

    \item \textbf{Unit\_Price and Quantity:}
    Selling price per unit and the
    quantity purchased.

    \item \textbf{Discount\_Percent:}
    Percentage discount applied
    to the transaction.

    \item \textbf{Cost\_Per\_Unit:}
    Cost associated with one unit
    of the product.

    \item \textbf{Marketing\_Spend:}
    Marketing expenditure associated
    with the transaction.

    \item \textbf{Customer\_Rating:}
    Rating recorded for the transaction.

    \item \textbf{Payment\_Method:}
    Method used to complete the payment.

    \item \textbf{Delivery\_Days:}
    Number of days required for delivery.

    \item \textbf{Returned:}
    Indicates whether the product
    was returned.

    \item \textbf{Sales\_Amount:}
    Total transaction sales value
    after applying the discount.

    \item \textbf{Profit:}
    Amount remaining after accounting
    for product costs and marketing
    expenditure.

\end{itemize}

The original dataset contained reported
data-quality issues, including missing
values and extreme numerical observations.

After preprocessing and feature engineering,
the dataset was expanded to
\textbf{58 columns}, including
29 one-hot encoded features and
8 scaled numerical features.

The processed dataset contains
no missing values or duplicate rows.

% =====================================
% 4.3 DATA TYPES
% =====================================

\section{Data Types}

The original dataset contains three
main categories of variables:
\textbf{numerical, categorical, and date}.

\textbf{1. Numerical Variables}

Numerical variables represent measurable
quantities and can be used in mathematical
and statistical calculations.

The dataset contains 10 numerical attributes:

\begin{itemize}
    \item Customer\_Age
    \item Unit\_Price
    \item Quantity
    \item Discount\_Percent
    \item Cost\_Per\_Unit
    \item Marketing\_Spend
    \item Customer\_Rating
    \item Delivery\_Days
    \item Sales\_Amount
    \item Profit
\end{itemize}

These variables are used for statistical
calculations, data analysis, and
dimensionality reduction.

\textbf{2. Categorical Variables}

Categorical variables represent labels,
identifiers, or classifications rather
than continuous numerical measurements.

The original dataset contains
10 categorical or identifier attributes:

\begin{itemize}
    \item Transaction\_ID
    \item Customer\_ID
    \item Customer\_Name
    \item Customer\_Type
    \item City
    \item Region
    \item Product\_ID
    \item Product\_Category
    \item Payment\_Method
    \item Returned
\end{itemize}

The \texttt{Returned} attribute is a
binary categorical variable represented
using numerical values.

Some categorical attributes were converted
into numerical representations using
one-hot encoding during feature engineering.

\textbf{3. Date Variable}

The \texttt{Date} attribute represents
the date associated with each transaction.

In the uploaded CSV file, this attribute
is stored as text. It can be converted
into a datetime data type for
date-based analysis.

\textbf{4. Engineered Features}

The processed dataset also contains
additional numerical features created
during preprocessing.

These include one-hot encoded columns
representing categorical values and
scaled columns representing standardized
numerical attributes.

The scaled features include customer age,
unit price, quantity, discount percentage,
cost per unit, marketing expenditure,
customer rating, and delivery duration.

These transformations prepare the dataset
for numerical analysis and techniques
such as Principal Component Analysis (PCA).

\endgroup

% =====================================
% 5. DATA PREPROCESSING
% =====================================


% =====================================
% CHAPTER 5: DATA PREPROCESSING
% =====================================

\chapter{Data Preprocessing}

\begingroup
\fontsize{16}{23}\selectfont

Data preprocessing is an important stage of
Data Science that involves identifying and
handling data-quality issues before applying
statistical and analytical techniques.

In the \textbf{Sales Insights} project,
preprocessing was performed to improve
the consistency and usability of the dataset.

The original dataset contained
\textbf{100 records and 21 attributes}.
During the initial data-quality assessment,
26 missing values were reported across
8 columns, affecting 22 records.

The preprocessing workflow involved examining
missing values, checking duplicate records,
identifying outliers, converting data types,
encoding categorical variables, and
standardizing numerical features.

% =====================================
% 5.1 MISSING VALUES
% =====================================

\section{Missing Values}

Missing values occur when certain observations
in a dataset do not contain recorded information.
These values can affect statistical calculations
and the performance of analytical algorithms.

In the original Sales Insights dataset,
missing values were reported in the
following attributes:

\begin{itemize}

    \item \texttt{Customer\_Age}
    \item \texttt{Customer\_Type}
    \item \texttt{City}
    \item \texttt{Discount\_Percent}
    \item \texttt{Marketing\_Spend}
    \item \texttt{Customer\_Rating}
    \item \texttt{Payment\_Method}
    \item \texttt{Delivery\_Days}

\end{itemize}

Missing numerical values can be handled
using statistical measures such as the
mean or median. The median is particularly
useful when the data contains extreme
observations because it is less affected
by outliers.

For categorical variables, mode imputation
can be used to replace missing values with
the most frequently occurring category.

In some cases, missing values may also
be estimated using known mathematical
relationships between existing attributes,
provided that these relationships are valid.

The processed dataset was checked for
remaining missing values using the
\texttt{isnull()} function in Pandas.

The final uploaded dataset contains
\textbf{zero missing values}.

% =====================================
% 5.2 DUPLICATE VALUES
% =====================================

\section{Duplicate Values}

Duplicate records occur when identical
observations appear more than once in
a dataset. Such records may affect
statistical summaries and introduce
unnecessary repetition into the analysis.

The Sales Insights dataset was examined
for duplicate rows using the
\texttt{duplicated()} function provided
by the Pandas library.

The dataset was also checked for
repeated transaction identifiers,
since each transaction should have
a unique \texttt{Transaction\_ID}.

No duplicate transactions were reported
in the original dataset.

The final processed dataset contains
\textbf{100 records with no duplicate rows}.

Therefore, no reduction in the number
of records was required due to
duplicate removal.

% =====================================
% 5.3 OUTLIER DETECTION
% =====================================

\section{Outlier Detection}


\begingroup
\fontsize{14}{21}\selectfont

In this project, the
\textbf{Interquartile Range (IQR) method}
was implemented to detect potential
outliers in the numerical attributes
of the dataset.

The first quartile ($Q_1$) and third
quartile ($Q_3$) were calculated for
the numerical variables, and the
Interquartile Range was determined
using their difference.

The lower and upper bounds were
then calculated using the
1.5 IQR rule. Observations falling
outside these bounds were identified
as potential outliers.

This method helped us identify
extreme numerical observations
and examine their distribution
within the dataset.

\endgroup

\chapter{Exploratory Data Analysis}

\begingroup
\fontsize{14}{21}\selectfont

Exploratory Data Analysis (EDA) is the process
of examining a dataset to understand its
characteristics, distributions, and relationships
between different variables.

In the \textbf{Sales Insights} project,
EDA was performed to examine the numerical
attributes of the dataset and understand
the distribution of sales transactions.

Descriptive statistics were calculated
using the \texttt{describe()} function
from the Pandas library. This provided
information such as the count, mean,
standard deviation, minimum, maximum,
and quartile values of numerical features.

Visualization techniques were also used
to represent the data graphically and
identify patterns that may not be
immediately visible from numerical values.

% =====================================
% VISUALIZATIONS
% =====================================

\section{Visualizations}

Data visualization is an important part
of Exploratory Data Analysis because
it helps represent numerical information
in a graphical format.

In this project, three visualization
techniques were used:
\textbf{Box Plot, Histogram, and
Correlation Heat Map}.

% =====================================
% BOX PLOT
% =====================================

\subsection{Box Plot}

A box plot is a statistical visualization
used to represent the distribution
and spread of numerical data.

It displays the median, first quartile
($Q_1$), third quartile ($Q_3$),
and the interquartile range (IQR).

The box represents the middle
50\% of the observations, while
the whiskers indicate the range
of values within the selected
outlier limits.

In the \textbf{Sales Insights} project,
a box plot was created for the
\texttt{Sales\_Amount} attribute.

The visualization was used to examine
the variation in sales amounts and
identify potential outliers or
extreme transaction values.

% =====================================
% HISTOGRAM
% =====================================

\subsection{Histogram}

A histogram is a graphical representation
of the frequency distribution of
a numerical variable.

It divides numerical observations
into intervals called bins and
displays the number of observations
falling within each interval.

In this project, a histogram was
created for the \texttt{Sales\_Amount}
attribute to examine the distribution
of sales values across transactions.

The histogram helped visualize
the frequency of different sales
ranges and observe whether
transactions were concentrated
around lower, middle, or
higher sales amounts.

It also provided a visual
representation of the overall
shape and spread of the
sales distribution.

% =====================================
% CORRELATION HEAT MAP
% =====================================

\subsection{Correlation Heat Map}

A correlation heat map is a
visualization technique used to
represent the relationships
between numerical variables.

It displays correlation coefficients
using different colours and
colour intensities.

The correlation coefficient
ranges from $-1$ to $+1$.

\begin{itemize}

    \item \textbf{Positive Correlation:}
    Indicates that two variables
    tend to increase together.

    \item \textbf{Negative Correlation:}
    Indicates that one variable
    tends to decrease as the
    other increases.

    \item \textbf{Near-Zero Correlation:}
    Indicates little or no
    linear relationship between
    two variables.

\end{itemize}

In the \textbf{Sales Insights}
project, a correlation heat map
was generated using the numerical
attributes of the dataset.

The visualization was used to
examine relationships between
different numerical features
and identify possible associations
between the variables.

The heat map provided a
graphical summary of the
correlation matrix and helped
in understanding the relationships
present in the dataset.

\endgroup


% =====================================
% FEATURE EXTRACTION AND ENGINEERING
% =====================================

\chapter{Feature Extraction and Feature Engineering}

\begingroup
\fontsize{14}{21}\selectfont

Feature extraction and feature engineering
involve creating or transforming variables
to obtain useful information from an
existing dataset.

In the \textbf{Sales Insights} project,
feature engineering was performed to
derive additional attributes from
the existing sales transaction data.

Three new features were created:
\textbf{Average Price Per Item},
\textbf{Profit Percentage}, and
\textbf{High Sales}.

% =====================================
% AVERAGE PRICE PER ITEM
% =====================================

\section{Average Price Per Item}

The feature \texttt{Avg\_Price\_Per\_Item}
was created by dividing the total
\texttt{Sales\_Amount} by the
\texttt{Quantity} purchased.

The formula used is:

\[
\text{Avg Price Per Item}
=
\frac{\text{Sales Amount}}{\text{Quantity}}
\]

This feature represents the average
sales value per unit in each transaction
after applying the discount.

It provides information about the
effective selling price of products
within individual transactions.

% =====================================
% PROFIT PERCENTAGE
% =====================================

\section{Profit Percentage}

The feature \texttt{Profit\_Percentage}
was calculated using the
\texttt{Profit} and
\texttt{Sales\_Amount} attributes.

The formula used is:

\[
\text{Profit Percentage}
=
\frac{\text{Profit}}{\text{Sales Amount}}
\times 100
\]

This feature represents the profit
earned as a percentage of the
total sales amount.

It helps in understanding the
relative profitability of each
transaction rather than considering
only the absolute profit value.

A higher profit percentage indicates
a greater proportion of profit
relative to the sales amount.

% =====================================
% HIGH SALES
% =====================================

\section{High Sales}

The feature \texttt{High\_Sales}
was created to classify transactions
into two categories based on
their sales amounts.

The median value of
\texttt{Sales\_Amount} was used
as the threshold for classification.

The feature was defined as:

\[
\text{High Sales} =
\begin{cases}
1, & \text{if Sales Amount} > \text{Median} \\
0, & \text{if Sales Amount} \leq \text{Median}
\end{cases}
\]

Transactions with sales amounts
greater than the median were
assigned a value of
\textbf{1 (High Sales)}.

Transactions with sales amounts
less than or equal to the median
were assigned a value of
\textbf{0 (Low Sales)}.

This feature provides a binary
representation of transaction
sales performance.

% =====================================
% SUMMARY
% =====================================

\section{Summary}

The three engineered features
provide additional information
about different aspects of
the sales transactions.

\begin{itemize}

    \item \textbf{Avg\_Price\_Per\_Item:}
    Represents the average sales
    value per unit.

    \item \textbf{Profit\_Percentage:}
    Represents the profit relative
    to the total sales amount.

    \item \textbf{High\_Sales:}
    Classifies transactions based
    on the median sales amount.

\end{itemize}

These features can support
further statistical analysis
and machine learning by providing
additional representations of
sales value, profitability,
and transaction performance.

\endgroup

\chapter{Hypothesis Testing}

\begingroup
\fontsize{14}{21}\selectfont

Hypothesis testing is a statistical method
used to determine whether an observed
difference between groups is statistically
significant.

In the \textbf{Sales Insights} project,
a Two-Sample Z-Test was performed to
compare the mean sales amounts of
Business and Retail customers.

The purpose of this test was to determine
whether the difference in average sales
between these two customer groups
is statistically significant.

% =====================================
% TWO-SAMPLE Z-TEST
% =====================================

\section{Two-Sample Z-Test}

The Two-Sample Z-Test is used to compare
the means of two independent groups
under suitable statistical assumptions.

In this project, the test was applied
to the \texttt{Sales\_Amount} variable
for Business and Retail customers.

% =====================================
% HYPOTHESES
% =====================================

\subsection{Hypotheses}

The following hypotheses were formulated:

\begin{itemize}

    \item \textbf{Null Hypothesis ($H_0$):}
    There is no difference between the
    population mean sales amounts of
    Business and Retail customers.

    \item \textbf{Alternative Hypothesis ($H_1$):}
    There is a difference between the
    population mean sales amounts of
    Business and Retail customers.

\end{itemize}

Mathematically,

\[
H_0: \mu_B = \mu_R
\]

\[
H_1: \mu_B \neq \mu_R
\]

where $\mu_B$ and $\mu_R$ represent
the population mean sales amounts
of Business and Retail customers,
respectively.

A significance level of
$\alpha = 0.05$ was selected.

% =====================================
% Z-TEST FORMULA
% =====================================

\subsection{Z-Test Formula}

The Z-statistic for comparing two
independent sample means can be
calculated using:

\[
Z =
\frac{\bar{X}_B-\bar{X}_R}
{\sqrt{
\frac{\sigma_B^2}{n_B}
+
\frac{\sigma_R^2}{n_R}
}}
\]

where:

\begin{itemize}

    \item $\bar{X}_B$ = Mean sales of
    Business customers.

    \item $\bar{X}_R$ = Mean sales of
    Retail customers.

    \item $\sigma_B^2$ = Population
    variance of Business customer sales.

    \item $\sigma_R^2$ = Population
    variance of Retail customer sales.

    \item $n_B$ = Number of Business
    customer observations.

    \item $n_R$ = Number of Retail
    customer observations.

\end{itemize}

When population variances are unknown,
sample variances may be used in a
large-sample Z approximation.

% =====================================
% TEST RESULTS
% =====================================

\subsection{Test Results}

The following results were obtained
from the Two-Sample Z-Test:

\begin{table}[H]
\centering
\renewcommand{\arraystretch}{1.4}

\begin{tabular}{|l|r|}
\hline
\textbf{Parameter} & \textbf{Value} \\
\hline
Business Customers & 28 \\
\hline
Mean Sales (Business) & INR 74,855.66 \\
\hline
Retail Customers & 67 \\
\hline
Mean Sales (Retail) & INR 37,617.20 \\
\hline
Z-Statistic & 2.944 \\
\hline
P-Value & 0.00324 \\
\hline
Significance Level & 0.05 \\
\hline
\end{tabular}

\caption{Two-Sample Z-Test Results}
\label{tab:ztest_results}

\end{table}

% =====================================
% INTERPRETATION
% =====================================

\subsection{Interpretation}

The calculated Z-statistic was
\textbf{2.944}, with a corresponding
two-sided p-value of
\textbf{0.00324}.

The significance level selected
for the test was
$\alpha = 0.05$.

Since,

\[
0.00324 < 0.05
\]

the p-value is smaller than
the significance level.

Therefore, the
\textbf{null hypothesis ($H_0$)
is rejected}.

This indicates that there is
a statistically significant
difference between the mean
sales amounts of Business and
Retail customers in the dataset.

The average sales amount for
Business customers was
\textbf{INR 74,855.66},
whereas the average sales amount
for Retail customers was
\textbf{INR 37,617.20}.

Thus, Business customers had
a higher average sales amount
than Retail customers in
the observed dataset.

The results suggest that
customer type is associated
with differences in average
transaction sales amounts.

However, statistical significance
does not establish a causal
relationship between customer
type and sales amount.

\endgroup



% =====================================
% MACHINE LEARNING IMPLEMENTATION
% =====================================

\chapter{Machine Learning Implementation}

\begingroup
\fontsize{14}{21}\selectfont

Machine Learning involves using data
and statistical algorithms to identify
patterns and make predictions.

In the \textbf{Sales Insights} project,
Simple Linear Regression was applied
to examine the relationship between
Marketing Spend and Sales Amount.

% =====================================
% LINEAR REGRESSION
% =====================================

\section{Linear Regression}

Simple Linear Regression is a supervised
machine learning technique used to
predict a continuous dependent variable
using one independent variable.

In this project, the following
variables were selected:

\begin{itemize}

    \item \textbf{Independent Variable (X):}
    Marketing Spend

    \item \textbf{Dependent Variable (Y):}
    Sales Amount

    \item \textbf{Number of Observations:}
    100

\end{itemize}

The objective was to examine whether
Marketing Spend could be used to
predict the Sales Amount of a transaction.

\medskip

\noindent\textbf{Mathematical Representation}

The equation of Simple Linear
Regression is:

\[
\hat{Y} = b_0 + b_1X
\]

where:

\begin{itemize}

    \item $\hat{Y}$ represents the
    predicted Sales Amount.

    \item $b_0$ represents the intercept.

    \item $b_1$ represents the slope
    of the regression line.

    \item $X$ represents Marketing Spend.

\end{itemize}

The slope and intercept were estimated
using the relationship between
Marketing Spend and Sales Amount
in the dataset.

% =====================================
% REGRESSION RESULTS
% =====================================

\subsection{Regression Results}

The following results were obtained
from the regression analysis:

\begin{table}[H]
\centering
\renewcommand{\arraystretch}{1.4}

\begin{tabular}{|l|r|}
\hline
\textbf{Parameter} & \textbf{Value} \\
\hline
Number of Observations & 100 \\
\hline
Mean Marketing Spend & INR 458.02 \\
\hline
Mean Sales Amount & INR 56,212.09 \\
\hline
Slope ($b_1$) & 28.9598 \\
\hline
Intercept ($b_0$) & 42,947.9957 \\
\hline
$R^2$ Score & 0.0118 \\
\hline
$R^2$ Percentage & 1.18\% \\
\hline
\end{tabular}

\caption{Simple Linear Regression Results}
\label{tab:linear_regression}

\end{table}

The regression equation obtained was:

\[
\boxed{
\hat{Y} = 42947.9957 + 28.9598X
}
\]

The positive slope indicates that
an increase of INR 1 in Marketing
Spend is associated with an increase
of approximately INR 28.96 in
predicted Sales Amount according
to the fitted regression model.

% =====================================
% MODEL INTERPRETATION
% =====================================

\subsection{Model Interpretation}

The performance of the regression
model was evaluated using the
Coefficient of Determination,
represented by $R^2$.

The obtained value was:

\[
R^2 = 0.0118
\]

Expressed as a percentage:

\[
R^2 \times 100 = 1.18\%
\]

This indicates that Marketing Spend
alone explains approximately
\textbf{1.18\% of the variation}
in Sales Amount within the dataset.

The low $R^2$ value indicates
that the fitted regression model
has very limited explanatory power.

Although the regression slope
is positive, the overall linear
relationship between Marketing
Spend and Sales Amount is weak.

Other variables, such as product
price, quantity, discount percentage,
and customer-related attributes,
may also help explain differences
in transaction sales amounts.

A scatter plot with the fitted
regression line can be used to
visualize the relationship between
the two variables.

The weak relationship is reflected
in the considerable variation of
observations around the regression
line.

\subsection{Conclusion}

The Simple Linear Regression analysis
demonstrated that Marketing Spend
alone is not a strong linear predictor
of Sales Amount in the given dataset.

The model explained only 1.18\%
of the observed variation in sales.

This analysis helped us understand
the implementation of Linear Regression,
the interpretation of regression
coefficients, and the importance
of evaluating model performance.

\endgroup


% =====================================
% DIMENSIONALITY REDUCTION
% =====================================

\chapter{Dimensionality Reduction}

\begingroup
\fontsize{14}{21}\selectfont

Dimensionality reduction is a technique
used to reduce the number of features
in a dataset while retaining important
information.

In the \textbf{Sales Insights} project,
dimensionality reduction was performed
using Principal Component Analysis (PCA).

The purpose was to simplify the numerical
data and represent it in a lower-dimensional
space for easier visualization and analysis.

% =====================================
% DIMENSIONALITY REDUCTION TECHNIQUES
% =====================================

\section{Dimensionality Reduction Techniques}

Dimensionality reduction techniques
transform high-dimensional data into
a smaller number of features.

These techniques are useful when
working with datasets containing
multiple numerical attributes.

In this project, Principal Component
Analysis was implemented to reduce
the selected numerical features
to two principal components.

% =====================================
% PRINCIPAL COMPONENT ANALYSIS
% =====================================

\subsection{Principal Component Analysis (PCA)}

Principal Component Analysis (PCA)
is an unsupervised dimensionality
reduction technique that transforms
the original numerical variables
into new variables called
principal components.

These components are linear
combinations of the original
features and are arranged
according to the amount of
variance they explain.

\medskip

\noindent\textbf{Data Standardization}

Before applying PCA, the following
10 numerical features were selected:

\begin{itemize}

    \item \texttt{Customer\_Age}
    \item \texttt{Unit\_Price}
    \item \texttt{Quantity}
    \item \texttt{Discount\_Percent}
    \item \texttt{Cost\_Per\_Unit}
    \item \texttt{Marketing\_Spend}
    \item \texttt{Customer\_Rating}
    \item \texttt{Delivery\_Days}
    \item \texttt{Sales\_Amount}
    \item \texttt{Profit}

\end{itemize}

These variables have different
numerical ranges and measurement
units. Therefore, standardization
was performed using
\textbf{StandardScaler} from
the Scikit-learn library.

The standardization formula is:

\[
z = \frac{x-\mu}{\sigma}
\]

where:

\begin{itemize}

    \item $x$ represents the
    original feature value.

    \item $\mu$ represents the
    mean of the feature.

    \item $\sigma$ represents the
    standard deviation.

    \item $z$ represents the
    standardized value.

\end{itemize}

Standardization helps prevent
features with larger numerical
scales from dominating PCA.

\medskip

\noindent\textbf{PCA Implementation}

After standardization, PCA was
applied to transform the
10 selected numerical features
into two principal components:

\begin{itemize}

    \item \textbf{PC1:}
    The first principal component,
    which captures the maximum
    possible variance in the
    standardized dataset.

    \item \textbf{PC2:}
    The second principal component,
    which captures the maximum
    remaining variance while being
    orthogonal to PC1.

\end{itemize}

The transformation reduced the
numerical feature space from
10 dimensions to 2 dimensions.

\medskip

\noindent\textbf{PCA Visualization}

A two-dimensional scatter plot
was created using PC1 and PC2.

Each point on the scatter plot
represents a sales transaction
from the dataset.

The plot was used to visualize
the distribution of transactions
in the reduced feature space
and examine possible patterns
or groupings.

\medskip

\noindent\textbf{Explained Variance Ratio}

The explained variance ratio
was calculated to determine
the proportion of total variance
captured by each principal component.

The formula is:

\[
EVR_i =
\frac{\lambda_i}
{\sum_{j=1}^{p}\lambda_j}
\]

where:

\begin{itemize}

    \item $\lambda_i$ represents
    the variance explained by
    the $i$th principal component.

    \item $p$ represents the
    total number of original
    numerical features.

\end{itemize}

The total explained variance
of the first two components is:

\[
EVR_{\text{total}} = EVR_1 + EVR_2
\]

This value represents the
proportion of total variance
retained after reducing the
dataset to two dimensions.

\medskip

\noindent\textbf{Conclusion}

Principal Component Analysis
was applied to reduce the
10 selected numerical features
into two principal components.

The resulting two-dimensional
representation was visualized
using a scatter plot.

The explained variance ratio
was used to assess how much
variance was retained by
the selected components.

This implementation helped
demonstrate the practical
application of PCA for
dimensionality reduction
and data visualization.

\endgroup

% =====================================
% MODEL EVALUATION
% =====================================

\chapter{Model Evaluation}

\begingroup
\fontsize{14}{21}\selectfont

Model evaluation was performed to assess
the effectiveness of the statistical and
machine learning techniques applied to
the \textbf{Sales Insights} dataset.

Different evaluation methods were used
depending on the type of analysis performed.
These included regression evaluation,
explained variance analysis, statistical
hypothesis testing, and visual examination
of the results.

% =====================================
% REGRESSION MODEL EVALUATION
% =====================================

\section{Regression Model Evaluation}

The Simple Linear Regression model was
evaluated using the Coefficient of
Determination ($R^2$), predicted values,
and residuals.

The model used \texttt{Marketing\_Spend}
as the independent variable and
\texttt{Sales\_Amount} as the
dependent variable.

The Coefficient of Determination
measures the proportion of variation
in the dependent variable explained
by the fitted regression model.

The obtained result was:

\[
R^2 = 0.0118
\]

Expressed as a percentage:

\[
R^2 \times 100 = 1.18\%
\]

This indicates that Marketing Spend
alone explains approximately
\textbf{1.18\% of the variation}
in Sales Amount.

The low $R^2$ value suggests
that the fitted regression model
has very limited explanatory power.

Predicted values and residuals
were also considered to understand
the difference between actual
and estimated sales amounts.

A residual is calculated as:

\[
e_i = y_i - \hat{y}_i
\]

where $y_i$ represents the actual
Sales Amount and $\hat{y}_i$
represents the predicted Sales Amount.

% =====================================
% PCA EVALUATION
% =====================================

\section{PCA Evaluation}

Principal Component Analysis (PCA)
was evaluated using the
\textbf{explained variance ratio}.

The explained variance ratio measures
the proportion of total variance
captured by each principal component.

In this project, PCA was used
to transform 10 selected numerical
features into two principal
components, PC1 and PC2.

The combined explained variance
of the two components is:

\[
EVR_{\text{total}} = EVR_1 + EVR_2
\]

The explained variance ratio
was used to assess how much
of the original standardized
data variation was retained
after dimensionality reduction.

A scatter plot of PC1 and PC2
was also created to visualize
the distribution of transactions
in the reduced two-dimensional
feature space.

% =====================================
% HYPOTHESIS TEST EVALUATION
% =====================================

\section{Hypothesis Test Evaluation}

A Two-Sample Z-Test was performed
to compare the mean Sales Amount
of Business and Retail customers.

The test was evaluated using
the Z-statistic and p-value.

The following results were obtained:

\begin{itemize}

    \item \textbf{Z-Statistic:} 2.944

    \item \textbf{P-Value:} 0.00324

    \item \textbf{Significance Level:}
    $\alpha = 0.05$

\end{itemize}

Since the p-value was smaller
than the significance level,

\[
0.00324 < 0.05
\]

the null hypothesis was rejected.

This indicates a statistically
significant difference between
the mean Sales Amount of Business
and Retail customers in the
given dataset, under the
assumptions of the test.

% =====================================
% VISUAL EVALUATION
% =====================================

\section{Visual Evaluation}

Data visualization techniques were
used to examine distributions,
relationships, and analytical
results within the dataset.

The following visualizations
were considered:

\begin{itemize}

    \item \textbf{Histogram:}
    Used to examine the frequency
    distribution of Sales Amount.

    \item \textbf{Box Plot:}
    Used to observe the spread
    of sales values and identify
    potential outliers.

    \item \textbf{Correlation Heat Map:}
    Used to examine linear
    relationships between
    numerical attributes.

    \item \textbf{PCA Scatter Plot:}
    Used to visualize the
    transformed dataset using
    PC1 and PC2.

    \item \textbf{Regression Plot:}
    Used to visualize the
    relationship between
    Marketing Spend and Sales Amount
    along with the fitted
    regression line.

\end{itemize}

These visualizations supported
the interpretation of statistical
results and helped identify
patterns in the dataset.

% =====================================
% OVERALL EVALUATION
% =====================================

\medskip

\noindent\textbf{Overall Evaluation}

The evaluation methods provided
different perspectives on the
techniques implemented in the
Sales Insights project.

The regression analysis showed
that Marketing Spend alone
explained only a small proportion
of the variation in Sales Amount.

The Z-test indicated a statistically
significant difference in mean
sales between Business and
Retail customers.

PCA provided a lower-dimensional
representation of the numerical
features, while the explained
variance ratio served as a
measure of information retention
in terms of variance.

Together, these evaluation
methods helped us interpret
the results of the statistical
and machine learning techniques
applied to the dataset.

\endgroup



% =====================================
% RESULTS AND FINDINGS
% =====================================

\chapter{Results and Findings}

\begingroup
\fontsize{14}{21}\selectfont

The analysis of the \textbf{Sales Insights}
dataset provided findings related to
sales distributions, customer segments,
feature relationships, dimensionality
reduction, and regression model performance.

The results obtained from Exploratory
Data Analysis (EDA), Feature Engineering,
Principal Component Analysis (PCA),
Hypothesis Testing, and Linear Regression
are summarized below.

% =====================================
% EDA FINDINGS
% =====================================

\section{EDA Findings}

Exploratory Data Analysis showed
considerable variation in numerical
attributes such as
\texttt{Sales\_Amount},
\texttt{Profit}, and
\texttt{Quantity}.

Some numerical variables showed
right-skewed distributions and
contained extreme observations.

The histogram and box plot
helped visualize the distribution
and spread of sales amounts.

The correlation heat map was
used to examine relationships
between numerical variables
and identify possible
linear associations.

These visualizations provided
a clearer understanding of
the characteristics of the dataset.

% =====================================
% FEATURE EXTRACTION FINDINGS
% =====================================

\section{Feature Extraction Findings}

Feature engineering was performed
to derive additional attributes
from the existing dataset.

The following features were created:

\begin{itemize}

    \item \textbf{Average Price Per Item:}
    Represents the average sales
    value per unit purchased.

    \item \textbf{Profit Percentage:}
    Represents profit as a percentage
    of the total sales amount.

    \item \textbf{High Sales:}
    Classifies transactions into
    high-sales and low-sales
    categories using the median
    Sales Amount.

\end{itemize}

These features provided additional
representations of transaction
value, profitability, and
sales performance.

% =====================================
% PCA FINDINGS
% =====================================

\section{PCA Findings}

Principal Component Analysis (PCA)
was applied to reduce the
dimensionality of the selected
numerical features.

The original 10 numerical features
were transformed into two
principal components,
\textbf{PC1 and PC2}.

This transformation provided
a two-dimensional representation
of the sales transactions.

A scatter plot of PC1 and PC2
was used to visualize the
distribution of observations
in the reduced feature space.

The explained variance ratio
was calculated to assess
the proportion of total
variance retained by
the principal components.

% =====================================
% HYPOTHESIS TESTING FINDINGS
% =====================================

\section{Hypothesis Testing Findings}

A Two-Sample Z-Test was performed
to compare the mean Sales Amount
of Business and Retail customers.

The following results were obtained:

\begin{itemize}

    \item \textbf{Business Customers:} 28

    \item \textbf{Mean Sales (Business):}
    INR 74,855.66

    \item \textbf{Retail Customers:} 67

    \item \textbf{Mean Sales (Retail):}
    INR 37,617.20

    \item \textbf{Z-Statistic:} 2.944

    \item \textbf{P-Value:} 0.00324

    \item \textbf{Significance Level:}
    $\alpha = 0.05$

\end{itemize}

Since the p-value was smaller
than the significance level,

\[
0.00324 < 0.05
\]

the null hypothesis was rejected.

The test indicated a
\textbf{statistically significant
difference} between the mean
sales amounts of Business
and Retail customers.

Business customers had a
higher average Sales Amount
than Retail customers
in the observed dataset.

% =====================================
% LINEAR REGRESSION FINDINGS
% =====================================

\section{Linear Regression Findings}

Simple Linear Regression was
applied to examine the
relationship between
\texttt{Marketing\_Spend}
and \texttt{Sales\_Amount}.

The fitted regression equation was:

\[
\hat{Y} = 42947.9957 + 28.9598X
\]

where $X$ represents Marketing
Spend and $\hat{Y}$ represents
the predicted Sales Amount.

The model produced an
$R^2$ score of:

\[
R^2 = 0.0118
\]

This corresponds to:

\[
1.18\%
\]

The result indicates that
Marketing Spend alone explained
approximately 1.18\% of the
variation in Sales Amount
within the dataset.

Therefore, the model showed
a very weak linear relationship
between Marketing Spend and
Sales Amount.

Marketing Spend alone was
not a strong predictor of
Sales Amount in the
fitted regression model.

% =====================================
% OVERALL FINDINGS
% =====================================

\section{Overall Findings}

The project demonstrated the
application of different
statistical and data science
techniques to the Sales
Insights dataset.

Exploratory Data Analysis
helped examine the distribution
and relationships of
numerical variables.

Feature engineering provided
additional attributes for
representing transaction
value and profitability.

Principal Component Analysis
reduced the selected numerical
features into two principal
components for visualization.

The Two-Sample Z-Test indicated
a statistically significant
difference in average sales
between Business and
Retail customers.

Linear Regression showed
that Marketing Spend alone
had limited ability to
explain the variation
in Sales Amount.

Overall, the findings suggest
that additional variables
may be needed to better
explain sales variation.

The project also provided
practical experience in
applying statistical analysis,
feature engineering,
dimensionality reduction,
and regression techniques
to a structured dataset.

\endgroup


% =====================================
% CONCLUSION
% =====================================

\chapter{Conclusion}

\begingroup
\fontsize{14}{21}\selectfont

The \textbf{Sales Insights} project provided
practical experience in applying various
data analysis, statistical, and machine
learning techniques to a structured
sales transaction dataset.

The project involved data preprocessing,
feature engineering, Exploratory Data
Analysis (EDA), Principal Component
Analysis (PCA), statistical hypothesis
testing using the Z-test, and
Simple Linear Regression.

Exploratory Data Analysis helped us
understand the distribution of numerical
variables, identify potential outliers,
and examine relationships between
different attributes using visualizations
such as histograms, box plots, and
correlation heat maps.

Feature engineering was used to create
additional variables, including
\textbf{Average Price Per Item},
\textbf{Profit Percentage}, and
\textbf{High Sales}. These features
provided alternative representations
of transaction values, profitability,
and sales performance.

Principal Component Analysis was
implemented to reduce 10 selected
numerical features into two principal
components, PC1 and PC2. This allowed
the dataset to be represented and
visualized in a two-dimensional space.

The \textbf{Two-Sample Z-Test} was used
to compare the average Sales Amount
of Business and Retail customers.
The obtained p-value of
\textbf{0.00324} was smaller than
the selected significance level
of 0.05.

Therefore, the null hypothesis
was rejected, indicating a
statistically significant difference
between the mean sales amounts
of the two customer groups,
under the assumptions of the test.

Simple Linear Regression was applied
to examine the relationship between
Marketing Spend and Sales Amount.
The model produced an
$R^2$ value of approximately
\textbf{0.0118 (1.18\%)}.

This indicates that Marketing Spend
alone explained only a small
proportion of the variation
in Sales Amount within the dataset.
Therefore, it was not a strong
linear predictor of sales
in the fitted model.

Overall, the project helped us
understand how different concepts
from the Fundamentals of Data Science
course can be implemented together
within a single analytical workflow.

As second-year B.Tech students
specializing in Data Science,
this project provided valuable
experience in working with datasets,
implementing statistical techniques,
interpreting results, and
collaborating as a team.

Since the dataset was synthetically
generated, the findings represent
patterns within the simulated data
and should not be interpreted
as conclusions about an actual
business or its customers.

\endgroup

\end{document}
